\section{sonification}

To overcome the limitations of passive music listening, the data must be transformed into an auditory format. This process relies on deliberate mapping strategies to ensure the physiological informations remains interpretable. According to Dubus and Bresin (2013), the effectiveness of such a system depends on how physical quantities are mapped onto acoustic dimensions like pitch, loudness, or timbre to eable accurate data perception \cite{dubus2013systematic}.

Dubus et al. (2013) performed a systematic review on sonification. They mention that sonification can enhance the perception of one's own body and is therefore widely used in rehabilitation and training. Among different sonification techniques, parameter mapping is especially promising because it allows for continuous transmission of information. In this process, a certain set of mappings between "physical dimensions" and "auditory dimensions" is defined \cite{dubus2013systematic}.

To establish these mappings, a three-stage process is required. First, the specific auditory dimensions to be used must be identified, such as pitch or timbre. Then the polarity has to be defined. This determines the direction in which the auditory dimension changes according to the data. For instance, if the heart rate increases, it must be decided whether the pitch should increase or decrease. Finally, the psychophysical scaling has to be addressed. Since humans perceive changes in sound non-linearly, it is necessary to define the magnitude of the auditory change relative to the data change \cite{dubus2013systematic} \cite{walker2000psychophysical}.

According to Dubus and Bresin (2013), many sonification approaches implement a strategy called "ecological perception". This means that certain sonification principles are part of our subconscious mind. This applies, for example, to loudness, which can be explained by the inverse distance law. A decrease in sound intensity is naturally associated with an increasing distance from a source. Such natural associations are deeply ingrained in human perception and should be prioritized when employing a sonification approach \cite{dubus2013systematic}.

To ensure a clear understanding of the mapping process, it is essential to define these auditory dimensions. Timbre refers to the perceived quality or 'color' of a sound, which is often used for the dimension 'Time'. In contrast, loudness relates to the sound's intensity or volume, making it the preferred mapping for 'Kinetics' to represent energy or force. Furthermore, pitch that is revered to as the perceived frequency of a sound, is commonly used to represent vertical data, while spatialization refers to the perceived location of a sound source in a 3D environment, often used to depict horizontal data \cite{dubus2013systematic}.

Furthermore, it is pointed out that the sonic material is decisive. If the sounds make the user uncomfortable, this can have drastic consequences for the effectiveness of the system. Therefore, the use of real-world sound samples is often preferred over simple synthetic tones \cite{dubus2013systematic}.

In addition, it is important to investigate the chosen sonification strategy scientifically. The study highlights that many approaches are currently chosen in an ad hoc or arbitrary manner without verifying their effectiveness through objective tests \cite{dubus2013systematic}.
